% ID: SEL-D0017-BB-E04
% ESTADO: validada. S018. Revisión matemática, didáctica y visual.
\textbf{a) Variación de la población.}\par\nopagebreak[4]
Como $e^{2t}>0$, se tiene $B'(t)>0$ durante las siete horas. La población es estrictamente creciente y nunca empieza a decrecer en ese intervalo.

\textbf{b) Primitiva y condición inicial.}\par\nopagebreak[4]
Al integrar se obtiene $B(t)=25000e^{2t}+K$. Como $B(0)=40000$, resulta $K=15000$ y
\[
\boxed{B(t)=25000e^{2t}+15000,\qquad 0\le t\le7.}
\]
Derivar la expresión y evaluar en cero recupera los dos datos del modelo.

\textbf{c) Valor a la hora y media.}\par\nopagebreak[4]
Sustituimos $t=1.5$ horas:
\[
B(1.5)=25000e^3+15000\approx517138.42.
\]
El modelo estima unas $\boxed{517138\text{ bacterias}}$.
